\documentclass[../../../include/open-logic-section]{subfiles}

\begin{document}

\olfileid{sth}{choice}{hartogs}
\olsection{د پرتلې وړتيا او د هارتوګس لم}

دا يې ښه اړخ و۔ اوس بد اړخ ته راځو۔ بې له انتخابه
چارې \emph{ګډې وډې} کېږي۔ د دې د ليدو دپاره د
\cite{Hartogs1915} دا ښکلې پايله وګورئ:

\begin{lem}[\emph{په $\ZF$ کښې}]\ollabel{HartogsLemma}
د هر سټ $A$ دپاره داسې ترتيبي عدد $\alpha$ شته
چې $\cardnless{\alpha}{A}$ وي۔
\end{lem}

\begin{proof}
که $B \subseteq A$ او $R \subseteq B^2$ وي، نو
\olref[ord-arithmetic][using-addition]{rankcomputation}
له مخې $\tuple{B, R} \subseteq
V_{\setrank{A}+4}$ وي۔ نو د بېلونې په
کارولو دا سټ په پام کښې ونيسئ:
\[
	C = \Setabs{\tuple{B, R} \in V_{\setrank{A}+5}}{B\subseteq A 
	\text{ او $\tuple{B, R}$ ښه ترتيب دے}}
\]
د تعويض او
\olref[ordinals][ordtype]{thmOrdinalRepresentation}
په کارولو دا سټ جوړ کړئ:
\[
	\alpha = \Setabs{\ordtype{B, R}}{\tuple{B, R} \in C}.
\]
د \olref[ordinals][basic]{corordtransitiveord} له مخې
$\alpha$ ترتيبي عدد دے، ځکه دا د ترتيبي عددونو انتقالي
سټ دے۔ ځکه که $\gamma \in \beta \in \alpha$ وي، نو
$\beta = \ordtype{B, R}$ د داسې کوم
$B \subseteq R$ دپاره وي، او بيا
$\gamma = \ordtype{B_b, R_b}$ د کوم $b
\in B$ دپاره وي، د
\olref[ordinals][iso]{wellordinitialsegment} له مخې؛
نو $\gamma \in \alpha$۔
\textbf{د سرچينې يادښت (OLSTH-023):} په تېر چاپي استدلال
کښې د فرعي سټ نسبت د مخکيني سټ پر ځاے له اړيکې سره
تړل شوے دے۔ دا د تعريف له شرط سره نه لګېږي۔ چاپي
رياضيکي نښه هماغسې ساتل شوې او دا تېروتنه نه ده رغول شوې۔

د تناقض دپاره فرض کړئ چې
!!a{injection} $f \colon \alpha \to A$ شته۔
بيا چې:
\begin{align*}
	B &= \ran{f}\\
	R &= \Setabs{\tuple{f(\alpha), f(\beta)} \in A \times A}{\alpha \in \beta}.
\end{align*}
څرګنده ده چې $\alpha = \ordtype{B, R}$ او
$\tuple{B, R} \in C$۔ نو $\alpha
\in \alpha$، چې تناقض دے۔
\end{proof}

له دې څخه يوه ژوره پايله راځي:

\begin{thm}[\emph{په $\ZF$ کښې}]
لاندې دعوې همارزې دي:
\begin{enumerate}
	\item\ollabel{equivwo} د ښه ترتيب اصل
	\item\ollabel{equivcompare} يا $\cardle{A}{B}$ وي يا
	$\cardle{B}{A}$ وي، د هر دوو سټونو $A$ او $B$
	دپاره
\end{enumerate}
\end{thm}

\begin{proof}
\emph{\olref{equivwo} $\Rightarrow$ \olref{equivcompare}.}
$A$ او $B$ وټاکئ۔ د \olref{equivwo} په کارولو
ښه ترتيبونه $\tuple{A, R}$ او $\tuple{B, S}$
شته۔ د
\olref[ordinals][ordtype]{thmOrdinalRepresentation}
له مخې $f \colon
\alpha \to \tuple{A, R}$ او $g \colon \beta \to \tuple{B, S}$
هم‌شکلي وګڼئ۔ د
\olref[sth][ordinals][basic]{ordinalsaresubsets} له مخې
يا $\alpha \subseteq \beta$ وي يا
$\beta \subseteq \alpha$ وي۔ که
$\alpha \subseteq \beta$ وي، نو
$\comp{f^{-1}}{g} \colon A \to B$ يوه
!!a{injection} ده، نو $\cardle{A}{B}$؛ او که
$\beta \subseteq \alpha$ وي، په همدې ډول
$\cardle{B}{A}$۔

\emph{\olref{equivcompare} $\Rightarrow$ \olref{equivwo}.}
$A$ وټاکئ؛ د \olref{HartogsLemma} له مخې داسې
ترتيبي عدد $\beta$ شته چې
$\cardnless{\beta}{A}$ وي۔ د
\olref{equivcompare} په کارولو
$\cardle{A}{\beta}$ لرو۔ نو يوه
!!{injection} $f \colon A \to
\beta$ شته، او د $A$ غړي د دې يو پر يو تابع په وسيله
د $\Setabs{\tuple{a, b} \in A \times A}{f(a)
\in f(b)}$ اړيکې په تعريفولو سره ښه ترتيبولے شو۔
\end{proof}
\noindent
سمدستي پايله دا ده: که ښه ترتيب ناکام شي، ځينې سټونه
د اندازې له مخې \emph{په رښتيا نه پرتله کېږي}۔ نو که
ښه ترتيب ناکام شي، د اصلي عددونو نامتناهيت هاخوا حساب
به ګډ وډ شي۔ د بېلګې په توګه، دا فکر به پرېږدو چې
که $A$ او $B$ نامتناهي وي، نو
$\cardeq{\cardeq{A \disjointsum B}{A \times
B}}{M}$، چې $M$ د $A$ او $B$ له منځه لوے
سټ دے (وګورئ
\olref[card-arithmetic][simp]{cardplustimesmax})۔
مسئله ساده ده: که د $A$ او $B$ اندازه
\emph{پرتله} نه شو کړے، نو دا پوښتنه بې‌معنا ده
چې کوم يو لوے دے۔

\end{document}